Applying formula 2.8, we obtain:
@@M0@@
Putting:
@@M1@@
we can also write this last relation as:
@@M2@@
Since the constant term of @@M3@@ is @@M4@@, and we are writing congruences modulo torsion, we can integrate and write the congruence @@M5@@. This gives the stated congruences.

\sgaStableTarget{sga6.e07.cor.4-4}\textbf{Corollary 4.4.} \uline{With the notation and hypotheses of} 4.3, \uline{we have for every} @@M6@@, \uline{and every} @@M7@@:
@@M8@@

The corollary follows immediately from 4.3, using the expression of the @@M9@@ as linear combinations with integer coefficients of the @@M10@@.

\sgaStableTarget{sga6.e07.rem.4-5}\textbf{Remark 4.5.} The proof of 4.3 itself suggests that (4.3.1) should hold not only modulo torsion, but also as an equality. In [RRR], the equality was proved in characteristic @@M11@@ by means of a geometric characterization of the @@M12@@. An extension of this method to all characteristics would probably be possible if one knew how to define exterior-power operations on perfect complexes (cf. XIV).
